\documentclass{article}
\usepackage{amsmath,amssymb,amsthm}
\begin{document}

\textbf{Subproblem 5 --- Obstruction from three edge directions.}

\medskip
\textbf{Setup.}
Let $S\subset\mathbb{R}^2$ be a finite set of $n\ge 3$ points in convex position,
let $T\subset\mathbb{R}^2$ be a finite set of translating vectors with signs
$\sigma(t)\in\{+1,-1\}$, both values present, and define
$\sigma(p)=\sum_{t\in T,\,p-t\in S}\sigma(t)$.

Assume the \textbf{extremal equality} $|\{p:\sigma(p)\neq 0\}|=n$.

The following facts from Subproblems 1--4 may be used without proof:
\begin{itemize}
  \item Vertices $s_1,\dots,s_n$ of $P=\operatorname{conv}(S)$ in cyclic order each have
        an associated unique translate $t_i\in T$ (the exposed translate on the normal cone).
  \item For each edge $[s_i,s_{i+1}]$ (with direction $a_i=s_{i+1}-s_i$), the set
        $F_i=T\cap[t_i,t_{i+1}]$ is an arithmetic progression
        $\{t_i,t_i+a_i,\dots,t_i+m_i a_i\}$ with alternating signs and $t_{i+1}=t_i+m_i a_i$.
\end{itemize}

Call two edges of $P$ \emph{parallel} if their direction vectors are positive scalar multiples
of each other (i.e.\ they point in the same direction along the boundary).  Say $P$ has
\emph{at least three edge-direction classes} if no two non-opposite boundary edges of $P$
are parallel (equivalently, $P$ is not a parallelogram and not degenerate in the sense
that every edge direction appears exactly once or in an antipodal pair with no third class).
More precisely: the edges of $P$ determine direction vectors $a_1,\dots,a_n$ (up to positive
scaling); $P$ has at least three edge-direction classes if the set
$\{[a_i]:1\le i\le n\}$ of direction classes has at least three elements.

\medskip
\textbf{Statement.}
If $P=\operatorname{conv}(S)$ has at least three distinct edge-direction classes, then
the extremal equality $|\{p:\sigma(p)\neq 0\}|=n$ cannot hold for any finite signed set $T$
with both signs present.

In particular:
\begin{itemize}
  \item Every triangle ($n=3$) has three distinct edge-direction classes, so $n=3$ is impossible.
  \item Every convex $n$-gon with $n\ge 5$ has at least three edge-direction classes
        (since a polygon cannot have opposite sides for all pairs when $n$ is odd, and for
        $n\ge 5$ even there is always a third direction), so all $n\ge 5$ are impossible.
  \item A non-parallelogram convex quadrilateral has three or four edge-direction classes,
        so it is also impossible.
\end{itemize}
Therefore the only surviving case is $n=4$ with exactly two edge-direction classes.

\medskip
\textbf{Hint.}
Suppose for contradiction that three distinct edge-direction classes $[a_i],[a_j],[a_k]$
exist.  By the alternating-chain structure, the sign of $t_i$ determines the signs of all
translates in $F_i$, and $\sigma(t_{i+1})=(-1)^{m_i}\sigma(t_i)$.
Follow the sign assignment around the boundary of $P$: starting at $t_1$ with sign
$\varepsilon_1=\sigma(t_1)$, after traversing all $n$ edges we return to $t_1$.
The sign after one full traversal is $(-1)^{m_1+m_2+\cdots+m_n}\varepsilon_1$.
For consistency this product must equal $\varepsilon_1$, so $\sum_i m_i$ must be even.
Now examine what the three distinct edge-direction classes force: two faces with the same
direction class (opposite edges) satisfy $a_j=\lambda a_i$ for some $\lambda>0$; but
the alternating chains on them produce translates that both start from the \emph{same}
vertex of $Q$ or overlap, creating either a repeated translate in $T$ (violating $T$ being
a set) or an extra cancellation-point not among the $n$ lonely vertices.
Formalize this by tracking the endpoint $t_{i+1}=t_i+m_i a_i$ around the full boundary
and showing that three direction classes force either a contradiction in $Q$'s geometry
or an extra non-canceling point.

\end{document}
